% Planned results: Q3-INFO0 4x768, MAE, 4 epochs, three training seeds per R0-R5. Screening layout; no final-scale R5 run or equivalence claim.
\begin{table}[t]
\centering\footnotesize
\caption{\tablelang{\textbf{Couplings R0--R5 and backbones} at the screening scale (MAE encoder, three training seeds each; Appendix~\ref{app:design_act}).}{\textbf{耦合方式 R0--R5 与主干}（筛选规模：MAE 编码器，各三个训练种子；附录~\ref{app:design_act}）。}}
\label{tab:supp_act}
\setlength{\tabcolsep}{4.4pt}
\begin{tabular}{@{}lccc@{}}
\toprule
\tablelang{Coupling}{耦合} & \tablelang{SR mean}{SR 均值} & \tablelang{SR std.}{SR 标准差} & p95 (ms) \\
\midrule
R0 \tablelang{(future losses off)}{（关闭未来损失）} & \resultplaceholder & \resultplaceholder & \resultplaceholder \\
R1 \tablelang{(queries, no edge)}{（查询，无连接）} & \resultplaceholder & \resultplaceholder & \resultplaceholder \\
R2 \tablelang{(world$\to$action)}{（世界$\to$动作）} & \resultplaceholder & \resultplaceholder & \resultplaceholder \\
R3 \tablelang{(action$\to$world)}{（动作$\to$世界）} & \resultplaceholder & \resultplaceholder & \resultplaceholder \\
R4 \tablelang{(bidirectional)}{（双向）} & \resultplaceholder & \resultplaceholder & \resultplaceholder \\
R5 \tablelang{(decoded geometry)}{（解码几何）} & \resultplaceholder & \resultplaceholder & \resultplaceholder \\
\bottomrule
\end{tabular}
\end{table}
